\documentclass[12pt]{article}

\usepackage{amsmath}
\usepackage{graphicx}
\usepackage{hyperref}
\usepackage{geometry}
\usepackage{array}
\geometry{margin=1in}
\usepackage{float}

\title{Spacecraft ADCS Homework 1 Report\\
NASA Starling Mission Analysis}
\author{Aleksander Garbuz \\ Israel Garcia}
\date{\today}

\begin{document}

\maketitle

% ==============================
% Changelog
% ==============================

\section*{Changelog}

\begin{center}
\begin{tabular}{| m{3cm} | m{4cm} | m{6cm} |}
\hline
\textbf{Date} & \textbf{Section Modified} & \textbf{Description of Changes} \\
\hline
\the\year-\the\month-\the\day & Initial Draft & Selected Starling mission. Added mission overview, ADCS architecture, mass estimate, and initial inertia model. \\
\hline
 & & \\
\hline
\end{tabular}
\end{center}

\vspace{1cm}

% ==============================
% Repository
% ==============================

\section*{Project Repository}

A GitHub repository has been created to store simulation code, numerical tools, and documentation for this project.

\textbf{Repository Link:}

\noindent
\url{https://github.com/ATG218/spacecraft-adcs-course}

All simulations are implemented in \textbf{Julia}, running in Jupyter notebooks (\texttt{.ipynb}).
Orbital dynamics are numerically integrated using the \texttt{DifferentialEquations.jl} package.
3D orbit animations are produced using \texttt{MeshCat.jl}.

\vspace{1cm}

% ==============================
% Section 1
% ==============================

\section{Course Project Logistics}

This homework begins the semester-long ADCS course project. 
For this project, we selected NASA’s \textbf{Starling} mission as the spacecraft platform for analysis.

The objectives of this project are to:

\begin{itemize}
    \item Model spacecraft rigid-body dynamics
    \item Analyze principal axis stability
    \item Characterize the attitude determination and control system
    \item Design and simulate control laws in later assignments
\end{itemize}

Each homework assignment will expand this document and ultimately be compiled into a final report.

% ==============================
% Section 2
% ==============================

\section{Mission Description and ADCS Requirements}

\subsection{Mission Overview}

NASA’s \textbf{Starling} mission consists of four identical 6U CubeSats operating in low Earth orbit. 
The mission demonstrates autonomous swarm technologies including:

\begin{itemize}
    \item Autonomous relative maneuver planning
    \item Optical-based relative navigation
    \item Inter-satellite networking
    \item Distributed onboard autonomy
\end{itemize}

The spacecraft were launched on July 17, 2023 and operate in a Sun-synchronous orbit at approximately 480 km altitude.

Although Starling is a multi-spacecraft mission, this project analyzes a single spacecraft as an isolated rigid body.

\subsection{Spacecraft Bus}

Each Starling spacecraft is built on a Blue Canyon Technologies (BCT) XB1 6U CubeSat bus. \cite{BCT_XB1}

The standard 6U CubeSat dimensions are approximately:

\[
0.10 \text{ m} \times 0.20 \text{ m} \times 0.30 \text{ m}
\]

The bus integrates:

\begin{itemize}
    \item Onboard avionics and computing
    \item Power generation and storage
    \item GNSS receiver
    \item Communications subsystem
    \item Attitude Determination and Control System (ADCS)
\end{itemize}

\subsection{Attitude Determination and Control System}

The Blue Canyon XB1 bus includes an integrated three-axis ADCS consistent with the Blue Canyon XACT family.

The system includes:

\begin{itemize}
    \item Star trackers for fine attitude determination
    \item Inertial measurement unit (gyroscopes)
    \item Reaction wheels for three-axis control
    \item Magnetorquers for momentum dumping
\end{itemize}

Blue Canyon 6U systems report pointing accuracy on the order of 0.002 degrees (1$\sigma$) for dual-star-tracker configurations.
The XACT family is rated for a slew rate of $\geq 10\,^\circ$/s for a 4~kg, 3U CubeSat
\cite{BCT_XACT}. Scaled to the Starling 6U bus mass of approximately 12~kg, the effective
slew rate is estimated at $3$--$5\,^\circ$/s, consistent with the higher inertia of the
larger spacecraft.

\subsection{ADCS Requirements}

Based on mission objectives and spacecraft architecture, the following capabilities are required \cite{SMAD}:

\begin{itemize}
    \item Three-axis stabilized control
    \item Reaction wheel authority to counter environmental torques
    \item Momentum dumping via magnetorquers
    \item Controlled slewing between pointing targets
\end{itemize}

Further quantitative analysis will follow in subsequent sections.


\section{Spacecraft Model}
\subsection{Body Reference Frame Definition}

A body-fixed reference frame $\mathcal{B} = \{\mathbf{b}_1, \mathbf{b}_2, \mathbf{b}_3\}$ 
is defined and attached to the spacecraft.
The origin of $\mathcal{B}$ is located at the spacecraft center of mass.

The axes of $\mathcal{B}$ are aligned with the principal geometric dimensions of the 
6U spacecraft bus:

\[
L_1 = 0.10 \text{ m}, \quad
L_2 = 0.20 \text{ m}, \quad
L_3 = 0.30 \text{ m}
\]

The body axes are defined as follows:

\begin{itemize}
    \item $\mathbf{b}_1$: aligned with the $Y$-direction, shortest dimension (0.10~m)
    \item $\mathbf{b}_2$: aligned with the $X$-direction, intermediate dimension (0.20~m)
    \item $\mathbf{b}_3$: aligned with the $Z$-direction, longest dimension (0.30~m)
\end{itemize}

The frame is defined to be right-handed such that
\[
\mathbf{b}_1 \times \mathbf{b}_2 = \mathbf{b}_3.
\]

The deployed solar panels are mounted on the $\pm \mathbf{b}_1$ faces of the spacecraft.
The hinge runs along the $\mathbf{b}_2$ direction (20~cm), and when deployed the panels
extend outward in the $\pm \mathbf{b}_1$ direction, lying approximately in the
$\mathbf{b}_1$--$\mathbf{b}_3$ plane.

This body-frame definition aligns the coordinate system with the dominant geometric
features of the spacecraft and simplifies subsequent inertia and surface discretization
calculations.

\subsection{Modeling, Center of Mass, and COM Inertia}

To gather the COM and Inertial Matrix for our chosen spacecraft, we first created a model in Fusion360


\subsection{Mass and Inertia Properties}

The exact mass of a Starling spacecraft is not publicly released. 
Comparable NASA 6U CubeSat missions report launch masses near 12 kg. 
We adopt 12 kg as a representative spacecraft mass.

Assuming a uniform-density rectangular prism with dimensions:

\[
x = 0.10 \text{ m}, \quad
y = 0.20 \text{ m}, \quad
z = 0.30 \text{ m}
\]

the principal moments of inertia are:

\[
I_x = \frac{1}{12} m (y^2 + z^2)
\]
\[
I_y = \frac{1}{12} m (x^2 + z^2)
\]
\[
I_z = \frac{1}{12} m (x^2 + y^2)
\]

Substituting \( m = 12 \) kg gives:

\[
I_x \approx 0.13 \text{ kg m}^2
\]
\[
I_y \approx 0.10 \text{ kg m}^2
\]
\[
I_z \approx 0.05 \text{ kg m}^2
\]

These values provide a first-order rigid-body inertia model.

% ==============================
% Section 4
% ==============================

\section{Dynamics}

\subsection{State Vector and Reference Frame}

The spacecraft position and velocity are expressed in an Earth-Centered Inertial (ECI) frame,
with units of kilometers and seconds throughout.
The 6-element state vector is
\[
\mathbf{x} = \begin{bmatrix} \mathbf{r} \\ \mathbf{v} \end{bmatrix} \in \mathbb{R}^6,
\]
where $\mathbf{r} = [x,\, y,\, z]^\top$~(km) is the ECI position and
$\mathbf{v} = [\dot{x},\, \dot{y},\, \dot{z}]^\top$~(km/s) is the ECI velocity.

\subsection{Equations of Motion}

The orbital ODE is
\[
\dot{\mathbf{r}} = \mathbf{v}, \qquad
\dot{\mathbf{v}} = \mathbf{a}_{2B} + \mathbf{a}_{J_2} + \mathbf{a}_{D}.
\]

\subsubsection*{Two-Body Gravity}

\[
\mathbf{a}_{2B} = -\frac{\mu}{r^3}\,\mathbf{r}, \qquad r = \|\mathbf{r}\|,
\]
with $\mu = 398600.4418$~km$^3$/s$^2$.

\subsubsection*{$J_2$ Oblateness Perturbation}

The first zonal harmonic of Earth's gravitational potential produces a secular drift in the
right ascension of the ascending node (RAAN).
In ECI coordinates the perturbation acceleration is \cite{Vallado2013}
\[
\mathbf{a}_{J_2} = \frac{3\,J_2\,\mu\,R_E^2}{2\,r^5}
\begin{bmatrix}
x \!\left(\dfrac{5z^2}{r^2} - 1\right) \\[8pt]
y \!\left(\dfrac{5z^2}{r^2} - 1\right) \\[8pt]
z \!\left(\dfrac{5z^2}{r^2} - 3\right)
\end{bmatrix},
\]
with $J_2 = 1.08262668 \times 10^{-3}$ and $R_E = 6378.1363$~km,
both consistent with the EGM96 gravity model \cite{EGM96}.

\subsubsection*{Atmospheric Drag}

An exponential atmosphere model is used with reference density
$\rho_0 = 3.614 \times 10^{-13}$~kg/m$^3$ at reference altitude $h_0 = 700$~km
and scale height $H = 88.667$~km \cite{Curtis2014}:
\[
\rho(h) = \rho_0 \exp\!\left(-\frac{h - h_0}{H}\right), \qquad h = \|\mathbf{r}\| - R_E.
\]
The drag acceleration is
\[
\mathbf{a}_D = -\frac{1}{2}\,\rho\,B_c\,|\mathbf{v}_{\mathrm{rel}}|\,\mathbf{v}_{\mathrm{rel}},
\]
where $\mathbf{v}_{\mathrm{rel}} = \mathbf{v} - \boldsymbol{\omega}_E \times \mathbf{r}$
accounts for Earth's rotation ($\omega_E = 7.2921159 \times 10^{-5}$~rad/s), and
$B_c = C_d A / m$ is the ballistic coefficient.
Representative values $C_d = 2.2$, $A = 0.03$~m$^2$, and $m = 12$~kg give
$B_c = 5.50 \times 10^{-3}$~m$^2$/kg.
These are approximate values; the cross-sectional area presented to the flow varies with
spacecraft attitude, which is not modeled in the current simulation.

\subsection{Sun-Synchronous Orbit Design}

A sun-synchronous orbit (SSO) requires the orbital plane to precess at Earth's mean
heliocentric rate of $\dot\Omega = +0.985647\,^\circ$/day.
The secular RAAN drift due to $J_2$ for a near-circular orbit is \cite{Vallado2013}
\[
\dot\Omega = -\frac{3}{2}\,J_2\,n\left(\frac{R_E}{p}\right)^2\!\cos i,
\qquad n = \sqrt{\frac{\mu}{a^3}},\quad p = a(1-e^2).
\]
Setting $\dot\Omega = +0.985647\,^\circ$/day and $e \approx 0$, the required inclination satisfies
\[
\cos i = \frac{\dot\Omega}{-\tfrac{3}{2}\,J_2\,n\,(R_E/a)^2}.
\]
For the Starling altitude of 480~km, $a = 6858.136$~km, and solving gives
$i = 97.326^\circ$, a retrograde inclination as expected for SSO.

\subsection{Initial Conditions}

Initial conditions are specified via classical orbital elements (COEs) and converted to ECI
coordinates using the standard perifocal-to-ECI direction cosine matrix
$Q = R_3(\Omega)\,R_1(i)\,R_3(\omega)$ \cite{Vallado2013}.
The epoch state is:

\begin{center}
\begin{tabular}{ll}
Semi-major axis   & $a = 6858.136$~km \\
Eccentricity      & $e = 0$ \\
Inclination       & $i = 97.326^\circ$ \\
RAAN              & $\Omega = 0^\circ$ \\
Arg.\ of perigee  & $\omega = 0^\circ$ \\
True anomaly      & $\nu = 0^\circ$ \\
\end{tabular}
\end{center}

This yields the initial ECI state
$\mathbf{r}_0 = [6858.136,\ 0,\ 0]^\top$~km and
$\mathbf{v}_0 = [0,\ {-0.9722},\ 7.5615]^\top$~km/s,
with orbital speed $|\mathbf{v}_0| = 7.6237$~km/s and
orbital period $T = 5652$~s (94.2~min).

\subsection{Numerical Simulation and Results}

The ODE is integrated using the Vern9 (9th-order Runge--Kutta) solver from the
Julia \texttt{DifferentialEquations.jl} package, with relative tolerance $10^{-10}$
and absolute tolerance $10^{-12}$.
The simulation spans 10 orbital periods ($\approx 15.7$~hours), with output saved every 30~s.
Two cases are run:

\begin{enumerate}
    \item \textbf{J2-only (no drag).}
    The RAAN precession rate extracted by a least-squares linear fit over 10 orbits is
    $\dot\Omega = 0.991\,^\circ$/day, compared to the target $0.985647\,^\circ$/day
    (error of $0.54\%$).
    The small discrepancy arises because the secular formula is a first-order orbital average,
    while the simulation propagates osculating elements that include short-period
    $J_2$ oscillations.
    The semi-major axis is conserved to within 43~m over 15.7~hours, confirming that
    $J_2$ produces no secular energy change for a circular orbit.

    \item \textbf{J2 + drag.}
    Over 15.7~hours, atmospheric drag reduces the semi-major axis by
    $\Delta a = -0.119$~km.
    In practice, Starling uses periodic maneuvers to maintain altitude; these are not
    modeled in the current simulation.
\end{enumerate}

Perturbations not included in the current model are solar radiation pressure, third-body
gravity, and spacecraft maneuvers; these are planned for future simulation work.
Figures~\ref{fig:orbit} and~\ref{fig:raan} show the simulated ECI trajectory
and RAAN evolution respectively.

\begin{figure}[H]
    \centering
    \includegraphics[width=0.75\textwidth]{figs/orbit_trajectory.png}
    \caption{Simulated ECI trajectory over 10 orbits (J2 + drag).
             Earth wireframe shown for scale.  Altitude 480~km, $i = 97.326^\circ$.}
    \label{fig:orbit}
\end{figure}

\begin{figure}[H]
    \centering
    \includegraphics[width=0.75\textwidth]{figs/raan_precession.png}
    \caption{RAAN evolution over 10 orbits (J2-only case).
             Simulated precession rate $0.991\,^\circ$/day versus SSO target $0.985647\,^\circ$/day.}
    \label{fig:raan}
\end{figure}

% ==============================
% References
% ==============================

\newpage

\begin{thebibliography}{9}

\bibitem{SMAD}
Wertz, J. R., and Larson, W. J.,
\textit{Space Mission Analysis and Design (SMAD)},
Microcosm Press.

\bibitem{STARLING_OVERVIEW}
NASA,
``Starling Mission Overview,''
NASA Mission Page,
\url{https://www.nasa.gov/mission/starling/}

\bibitem{STARLING_FLIGHT}
Miller, S. et al.,
``Starling CubeSat Swarm Technology Demonstration Flight Results,''
NASA Technical Reports Server (NTRS), 2024,
\url{https://ntrs.nasa.gov/}

\bibitem{BCT_XB1}
Blue Canyon Technologies,
``XB1 6U CubeSat Bus Overview,''
Product Documentation,
\url{https://www.bluecanyontech.com/spacecraft/6u-cubesat/}

\bibitem{BCT_XACT}
Blue Canyon Technologies,
``XACT Attitude Determination and Control System,''
Product Documentation,
\url{https://www.bluecanyontech.com/spacecraft-subsystems/xact/}

\bibitem{Vallado2013}
D.~A. Vallado,
\textit{Fundamentals of Astrodynamics and Applications}, 4th~ed.,
Microcosm Press / Springer, 2013.

\bibitem{Curtis2014}
H.~D. Curtis,
\textit{Orbital Mechanics for Engineering Students}, 3rd~ed.,
Butterworth-Heinemann / Elsevier, 2014.

\bibitem{EGM96}
F.~G. Lemoine et al.,
``Development of the Joint NASA GSFC and NIMA Geopotential Model EGM96,''
NASA/TP-1998-206861, NASA Goddard Space Flight Center, 1998.

\end{thebibliography}

\end{document}